\section{Ambiguity: de Pairwise Message a Hough-Style Voting}

La evidencia de bajo nivel es ambigua por naturaleza: un circle puede ser left eye, right eye, front wheel o back wheel. Si cada part candidata envía directamente a todas las demás parts candidatas un pose-transformed message del tipo «¿me apoyas?», los relation types y los spatial pairs crecen con rapidez. GLOM propone otra dirección: que cada part prediga la identity-pose de su parent y que después, en el parent level, se verifique si esas predicciones coinciden.

\lecturefigure{slide-24-upper-level-disambiguation.jpg}{Uso de la spatial relation correcta para resolver ambiguous parts en el parent level.}{Video oficial de Stanford Online 00:39:04--00:43:35.}

\subsection{La complejidad del transformational random field}

Supongamos que la nose hypothesis necesita preguntar si cerca existe una mouth compatible. El message no dice solo «hay / no hay mouth»: también tiene que transformar la nose pose mediante la nose-to-mouth relation en una expected mouth pose y, de regreso, aplicar la inverse transform. Si $N$ posiciones candidatas y $H$ tipos de relation head interactúan directamente, el número de messages en el esquema ingenuo se acerca a $O(HN^2)$, y cada message incluye una coordinate transform.

\begin{warningbox}{Aquí $O(N^2)$ es una presión de diseño, no datos medidos con un profiler}
En la clase se usa para explicar por qué el direct relation-specific messaging es complejo. El costo real depende además de la local window, el candidate pruning, el head sharing, la sparsity y la representation dimension; no puede usarse para clasificar directamente por velocidad a GLOM frente a la Hough route.
\end{warningbox}

\subsection{Hough transform: que las parts converjan en el parent space}

Una alternativa más sencilla es que nose y mouth no se reconozcan entre sí de forma directa, sino que cada una prediga la identity y la pose de la face. Si dos parts con posiciones y apariencias distintas producen el mismo parent embedding, la attention del parent level las considerará como apoyo mutuo. Esto se parece al Hough voting clásico: la evidence local se proyecta en un hypothesis space común, y los picos representan explicaciones consistentes.

\lecturefigure{slide-25-hough-transform.jpg}{Ruta Hough-style: las parts predicen el whole y después se compara si las whole hypotheses coinciden.}{Video oficial de Stanford Online 00:43:36--00:44:31.}

\begin{knowledgebox}{Lectura de la figura: comparación entre direct relation y common-parent vote}
La direct route, a la izquierda de la figura, requiere que la nose envíe relation-specific messages a distintas mouths, eyes, etc. potenciales; la parte derecha proyecta cada part en el whole hypothesis space. Aunque nose y mouth están en columns distintas, basta con que predigan el mismo face identity-pose vector para apoyarse mutuamente en el face level. La comparación decisiva no es «si hay o no attention», sino si la attention ocurre en el pairwise part space o en el common parent space.
\end{knowledgebox}

La parent prediction de una part $p$ puede abstraerse como:

\[
\hat{\mathbf e}_{w}^{(p)}=f_{\mathrm{bu}}\!\left(\mathbf e_p\right).
\]

Si nose y mouth pertenecen a la misma face, el caso ideal es

\[
\hat{\mathbf e}_{\mathrm{face}}^{(\mathrm{nose})}
\approx
\hat{\mathbf e}_{\mathrm{face}}^{(\mathrm{mouth})}.
\]

Aquí, $\mathbf e_p$ contiene la part identity y la pose, $f_{\mathrm{bu}}$ aprende la part-to-whole relation y $\hat{\mathbf e}_w^{(p)}$ es el vote de esa part por su parent. La agreement respalda un parent común, pero la disagreement también puede deberse a occlusion, a una bad pose estimate o a varios objects superpuestos.

\begin{importantbox}{Por qué esta ruta reduce el routing burden}
El embedding de cada column describe siempre «qué es, en cierto nivel» esa posición espacial; dentro de la column no hace falta mover dinámicamente la activity a otra capsule dedicada. Las conexiones entre columns siguen existiendo, pero se reducen a same-level similarity attention: lo que se compara son los parent votes, no una routing table dedicada para cada part-pair.
\end{importantbox}

\subsection{Cómo conservar multiple hypotheses}

En la early inference, un circle no puede reducirse de inmediato a una única explicación. Hinton propone que cada neuron del embedding aporte una log-probability basis function amplia sobre el joint identity-pose space $z=(\text{identity},\text{pose})$; al sumar las contribuciones de varias active neurons, puede formarse una distribución aguda o multimodal.

\lecturefigure{slide-26-joint-identity-pose.jpg}{Representación de una multimodal prediction mediante una joint identity-pose log-probability basis.}{Video oficial de Stanford Online 00:44:32--00:44:45.}

\begin{knowledgebox}{Lectura de la figura: solo al sumar vague bases se forma una sharp hypothesis}
Una sola neuron aporta una log-probability surface amplia sobre el joint identity-pose space, por lo que no puede leerse por sí sola como «esta neuron es el ojo izquierdo». La suma de varias active bases equivale a la intersección de varias constraints y puede dejar uno o varios picos. La figura no muestra la normalized probability, el método de entrenamiento ni la calibration; solo explica cómo la distributed activity podría, en principio, conservar multiple hypotheses.
\end{knowledgebox}

Una formulación didáctica es:

\[
\log p(z\mid\mathbf e)=\sum_{i=1}^{d} e_i\,\phi_i(z)-\log Z(\mathbf e).
\]

Aquí, $z$ es la identity-pose hypothesis, $e_i$ es la $i$-ésima activity del embedding, $\phi_i(z)$ es la basis function amplia de esa neuron en el hypothesis space y $Z(\mathbf e)$ es la normalization constant. Sumar las log probabilities de distintas evidences equivale a multiplicar las probabilities, lo que permite que las regiones de apoyo superpuestas se vuelvan más agudas.

\begin{warningbox}{El propio expositor llama a este razonamiento un weak argument}
Se trata de una representational hypothesis: explica cómo la high-dimensional distributed activity \emph{podría} expresar multimodal uncertainty. En la clase no se presentan calibrated likelihood, basis visualization, ablation ni neural recording que validen la fórmula. Las notas conservan la fórmula para que la hipótesis pueda ponerse a prueba, no para darle una falsa autoridad empírica.
\end{warningbox}

\teachervoice{Hinton dice que la perception debe manejar uncertainty, por lo que una neuron no puede representar un solo objeto determinado; pero también reconoce que «es la única forma que se me ocurre» no constituye una prueba sólida. Esta autocrítica debe estar en la misma página que la fórmula; de lo contrario, el lector tomará una preferencia de diseño por un hecho.}

\subsection{Resumen de la sección}

GLOM usa el part-to-whole vote para convertir el relation-specific pairwise messaging en parent-level agreement; después usa una distributed log-probability basis para dejar espacio a la early ambiguity. Ambos pasos tienen mucho poder explicativo, pero requieren, respectivamente, un complexity experiment y una uncertainty calibration para convertirse en mecanismos validados.
